\chapter{Background and Literature Review}
\label{ch:background}

This chapter first sets out the technical background needed to follow the rest of
the report, then reviews the research literature and the existing applications
closest to LectureAssist, and finally identifies the gaps that this project sets
out to fill.

\section{Preliminaries}
\label{sec:back-prelim}

\subsection{Large Language Models and Prompting}
\label{sec:back-llm}

A large language model (LLM) is an autoregressive neural network trained to predict
the next token of a sequence; at inference it generates text by sampling from the
distribution it has learned. Modern instruction-tuned LLMs can be steered entirely
through the prompt, without any gradient update, which is what makes them practical
for a student project: the same frozen model can be turned into a summariser, a
question writer, a grader or a judge purely by changing the text it is given. This
project uses the open-weight Gemma~3 family~\cite{gemma3}, served locally by the
Ollama runtime~\cite{ollama}, in two sizes: a fast $4$-billion-parameter model for
generation, and a $12$-billion-parameter model for the tasks where quality matters
more than latency (summarisation, grading, judging and answer verification).

Two prompting strategies are directly relevant, because they are used in this work
both as baselines and as generation methods in their own right.

\paragraph{Chain-of-Thought (CoT).} Chain-of-thought prompting asks the model to
produce its intermediate reasoning in natural language \emph{before} stating a final
answer, rather than jumping straight to it~\cite{wei2022chain}. Making the reasoning
explicit improves accuracy on multi-step problems, because each intermediate step
conditions the next and the model is no longer forced to commit to an answer in a
single forward pass.

\paragraph{Program-of-Thought (PoT).} Program-of-thought is a variant in which the
reasoning is expressed as an explicit, ordered sequence of \emph{computational}
steps --- in the original formulation, as executable code --- rather than as
free-form prose, so that the answer is \emph{computed} rather than
recalled~\cite{chen2023program}. It is particularly effective when the answer is a
numeric quantity, which is exactly the situation in a mathematics chapter, because
it separates the act of reasoning about \emph{what} to compute from the act of
performing the arithmetic.

\subsection{Retrieval-Augmented Generation}
\label{sec:back-rag}

An LLM only knows what is in its weights and in its context window. Retrieval-augmented
generation (RAG)~\cite{lewis2020rag} closes that gap: the source material is split
into chunks, each chunk is embedded into a dense vector, and at query time the
chunks whose embeddings are most similar to the query are retrieved and pasted into
the prompt, so the model answers \emph{from the document} rather than from memory.
LectureAssist embeds lecture chunks with the Sentence-BERT model
\texttt{all-MiniLM-L6-v2}~\cite{reimers2019sbert} and indexes them with
FAISS~\cite{johnson2019faiss} using inner-product search over $L_2$-normalised
vectors, which makes the similarity score the cosine similarity

\begin{equation}
\label{eq:cosine}
\mathrm{sim}(q, c) \;=\; \frac{\mathbf{e}_q \cdot \mathbf{e}_c}
                              {\lVert \mathbf{e}_q \rVert \, \lVert \mathbf{e}_c \rVert} .
\end{equation}

The same embedding machinery is reused at evaluation time, where it supplies two
LLM-free metrics: how well a generated question is \emph{grounded} in the lecture
(its best cosine match against any lecture chunk), and how \emph{diverse} a question
set is (one minus the mean nearest-neighbour cosine between its questions).

\subsection{Agentic Architectures}
\label{sec:back-agentic}

An \emph{agentic} system does not treat the LLM as a single-shot text generator. It
gives the model a goal, a set of tools it may call, and a loop in which it can act,
observe the result, and act again. The ReAct pattern~\cite{yao2023react} interleaves
reasoning traces with tool calls; reflection-style methods such as
Reflexion~\cite{shinn2023reflexion} add a self-critique step whose verbal feedback
is fed back into the next attempt. LectureAssist applies both ideas to question
generation: the generator may call a retrieval tool when it judges its context
insufficient, and a separate validator agent critiques each draft question and
selects a remediation action that determines how the next attempt differs from the
last.

\subsection{Speech Recognition and Optical Character Recognition}
\label{sec:back-asr-ocr}

To generate questions from a lecture \emph{video}, the lecture must first be turned
into text through two independent channels. Whisper~\cite{radford2023whisper} is a
transformer-based automatic speech recognition model trained on large-scale weakly
supervised data, robust to accents and background noise; this project uses the
Faster-Whisper implementation of the \texttt{small} checkpoint for the audio
channel. The visual channel is handled by optical character recognition, which
extracts text from sampled video frames. OCR accuracy is highly sensitive to
preprocessing, and the correct preprocessing depends on what is on screen --- white
chalk on a dark board needs to be inverted, a bright projected slide does not ---
which motivates the scene-adaptive design described in
Chapter~\ref{ch:methodology}.

\subsection{LLM-as-a-Judge}
\label{sec:back-judge}

Evaluating hundreds of generated questions by hand is not feasible within a
capstone project, and $n$-gram overlap metrics such as BLEU are known to correlate
poorly with question quality --- a question can be excellent and share almost no
vocabulary with any reference. The now-standard alternative is
\emph{LLM-as-a-judge}~\cite{zheng2023judging}: a strong model is given a rubric and
scores each item. It is a useful but imperfect instrument --- judges are known to
favour fluent and verbose text, and a judge from the same model family as the
generator will share its blind spots --- so this project treats the judge as a
\emph{relative} signal between pipelines, deliberately uses a larger model as the
judge than as the generator, and cross-checks the judge against a solver calibrated
on questions with known answers.

\section{Literature Review}
\label{sec:back-lit}

\subsection{Related Research}
\label{sec:back-related}

Automatic question generation (AQG) has a long history, surveyed comprehensively by
Kurdi et al.~\cite{kurdi2020systematic}. The earliest systems were
\emph{rule- and template-based}: they applied syntactic transformations to a
declarative sentence to turn it into an interrogative one, then ranked the
candidates with a supervised model, as in the work of Heilman and
Smith~\cite{heilman2010good}. These systems are transparent and grammatical, but
they are fundamentally limited to questions whose answer is a span of a single
sentence, and they cannot ask anything that requires synthesis across a passage.

The neural era reframed the task as sequence-to-sequence learning. Du et
al.~\cite{du2017learning} trained an attention-based encoder--decoder to generate
questions directly from a sentence or paragraph, outperforming rule-based systems on
human evaluation, and subsequent work fine-tuned pretrained transformers on
SQuAD-style corpora. Two limitations persisted. First, these models were trained on
extractive reading-comprehension data, so they inherited its bias toward shallow,
factoid questions. Second, they generate a question and an answer span but have no
mechanism to \emph{check} either one.

The current wave prompts general-purpose LLMs to write questions, including
multiple-choice questions with distractors, which sidesteps the need for
task-specific training data entirely and produces markedly more fluent items. The
open problem has shifted accordingly: fluency is no longer the bottleneck,
\emph{correctness and diversity} are. An LLM will reliably produce a well-formed MCQ
whose marked option is wrong, or whose four options are four paraphrases of the same
idea, and it will do so with complete confidence. Evaluations in this line of work
overwhelmingly report quality-style ratings --- fluency, relevance, answerability
--- and rarely verify that the marked answer is the right one.

Orthogonally, the prompting literature has shown that \emph{how} the model is asked
matters as much as which model is asked. Chain-of-thought~\cite{wei2022chain} and
program-of-thought~\cite{chen2023program} prompting substantially improve LLM
accuracy on multi-step and numerical reasoning, and agentic
patterns~\cite{yao2023react, shinn2023reflexion} improve it further by letting the
model retrieve evidence and revise its own output. These techniques are well
established for \emph{answering} questions. Their value for \emph{generating} them
--- where the model must derive an answer it then has to mark correctly --- has
received far less attention, and is one of the questions this project sets out to
answer empirically.

\subsection{Similar Applications}
\label{sec:back-apps}

Several commercial and open products overlap with parts of LectureAssist.
Flashcard and quiz platforms such as Quizlet and Anki support spaced retrieval
practice, but the questions must be authored by a human; they solve the
\emph{practice} problem and not the \emph{supply} problem. Note-taking assistants
such as Notion~AI and NotebookLM will summarise an uploaded document and answer
questions about it, and NotebookLM in particular is grounded in the user's own
sources, but their assessment support is shallow and they are cloud services.
Lecture-capture and transcription tools such as Otter.ai produce a transcript from
audio but stop there: they do not read the board, and they generate no assessment.
General-purpose chatbots such as ChatGPT will happily write quiz questions from
pasted text, but they require the student to supply clean text, they perform no
validation of the answer they mark, they retain no per-topic model of the student
across sessions, and they are neither free nor local.

Two properties are consistently absent across all of them. None of these tools reads
the \emph{board} --- the derivation the lecturer actually wrote out is invisible to
a transcript-only pipeline --- and none of them runs on the student's own hardware
without a paid API.

\section{Gap Analysis}
\label{sec:back-gap}

The review above identifies four gaps, which this project addresses directly.

\begin{enumerate}
  \item \textbf{The pipeline is not closed.} Existing work operates on
        pre-extracted, clean text. The step that is actually hardest for a student
        --- getting the lecture out of a video, including what was written on the
        board --- is assumed away. LectureAssist closes the loop from raw media to
        graded assessment, treating the visual and audio channels as first-class
        and complementary inputs.

  \item \textbf{Generation is unvalidated.} Neural and LLM-based generators emit a
        question and mark an answer, with no mechanism to check that the marked
        answer is correct or that the question is answerable from the source.
        LectureAssist adds a validator agent inside the generation loop, and, at
        evaluation time, a \emph{separate} Answer Generator agent that re-solves
        every question without seeing the marked option, so that answer correctness
        is measured rather than assumed.

  \item \textbf{Evaluation measures the wrong thing.} The field reports fluency-style
        quality scores, which LLM judges saturate. This project reports answer
        correctness and duplication alongside quality, and it combines them into a
        duplicate-penalised accept rate --- a decision that turns out to matter,
        because it reverses the ranking of the pipelines at the hard difficulty
        level (Chapter~\ref{ch:results}).

  \item \textbf{Agentic and prompting strategies are untested for generation.} CoT,
        PoT and agentic loops are well studied for \emph{answering} questions but
        rarely compared head-to-head for \emph{generating} them. This project runs
        that comparison on identical content, with an identical judge, at three
        difficulty levels.
\end{enumerate}

Finally, all of the above must hold under a hard deployment constraint: no paid API
and a single consumer-grade GPU, so that the system is usable in exactly the
bandwidth- and cost-limited settings where the need is greatest.
